On considère la fonction $f$ définie pour tout réel $x$ par :

\[ f(x) = \ln \left(\e^{\frac{x}{2}} + 2\right) \]

On admet que la fonction $f$ est dérivable sur $\mathbb{R}$.

On considère la suite $(u_{n})$ définie par $u_{0} = \ln (9)$ et, pour tout entier naturel $n$,

\[ u_{n+1} = f(u_{n}) \]

\smallskip

\begin{enumerate}
\item Montrer que la fonction $f$ est strictement croissante sur $\mathbb{R}$.
\item Montrer que $f(2\ln (2)) = 2\ln (2)$.
\item Montrer que $u_{1} = \ln (5)$.
\item Montrer par récurrence que pour tout entier naturel $n$, on a :

\[2\ln(2) \leqslant u_{n+1} \leqslant u_{n}\]

\item En déduire que la suite $(u_{n})$ converge.

\item
	\begin{enumerate}
		\item Résoudre dans $\R$ l'équation $X^{2} - X - 2 = 0$.
		\item En déduire l'ensemble des solutions sur $R$ de l'équation :

\[\e^{x} - \e^{\frac{x}{2}} - 2 = 0\]

		\item En déduire l'ensemble des solutions sur $\R$ de l'équation $f(x) = x$.
		\item Déterminer la limite de la suite $(u_{n})$.
	\end{enumerate}
\end{enumerate}
